% titel: Blatt 2, A6: Empirische Formel aus Verbrennungsanalyse
% Teil II: der durchgerechnete Weg.

\subsection*{Gegeben und gesucht}
Probe $m = \qty{255}{mg}$ (nur C, H, ggf. O). Verbrennungsprodukte:
$m_{\mathrm{CO_2}} = \qty{561}{mg}$, $m_{\mathrm{H_2O}} = \qty{306}{mg}$.
Gesucht: empirische Formel.

\subsection*{Stoffmengen C und H (alle Werte in Gramm umrechnen)}
\begin{align*}
  n(\mathrm{C}) &= \frac{0.561}{44.011} = 1.275\cdot10^{-2}\ \mathrm{mol}
    = \qty{12.75}{mmol},\\
  n(\mathrm{H}) &= 2\cdot\frac{0.306}{18.015}
    = 2 \cdot 1.699\cdot10^{-2} = 3.397\cdot10^{-2}\ \mathrm{mol}
    = \qty{33.97}{mmol}.
\end{align*}

\subsection*{Massen von C und H in der Probe}
\begin{align*}
  m(\mathrm{C}) &= 1.275\cdot10^{-2} \cdot 12.011 = \qty{153.1}{mg},\\
  m(\mathrm{H}) &= 3.397\cdot10^{-2} \cdot 1.008 = \qty{34.2}{mg}.
\end{align*}

\subsection*{Sauerstoffmasse und Stoffmenge}
\begin{equation}
  m(\mathrm{O}) = 255 - 153.1 - 34.2 = \qty{67.7}{mg}
  \qquad\Rightarrow\qquad
  n(\mathrm{O}) = \frac{0.0677}{15.999} = \qty{4.23}{mmol}.
\end{equation}

\subsection*{Verhältnis}
\begin{equation}
  n(\mathrm{C}):n(\mathrm{H}):n(\mathrm{O})
  = 12.75:33.97:4.23 \ \Big/\, 4.23 \ \Rightarrow\ 
  \frac{n(\mathrm{C})}{n(\mathrm{O})}\approx 3,01,\quad
  \frac{n(\mathrm{H})}{n(\mathrm{O})}\approx 8,03.
\end{equation}

\begin{ergebnis}
\begin{equation}
  \text{empirische Formel: }\ \mathrm{C_3H_8O}.
\end{equation}
\end{ergebnis}

\begin{kontrolle}
Verhältnisse liegen nahe den ganzen Zahlen $3:8:1$. Die Masseprobe passt:
$m(\mathrm{C}) + m(\mathrm{H}) + m(\mathrm{O}) = 153.1 + 34.2 + 67.7 =
255.0\ \mathrm{mg}$, exakt die eingesetzte Probemasse.
\end{kontrolle}

\begin{merksatz}
Aus den Verbrennungsprodukten $n(\mathrm{C})$ und $n(\mathrm{H})$ bestimmen
(Faktor 2 beim Wasserstoff), Sauerstoff als Differenz zur Probemasse ermitteln,
dann ins kleinste ganzzahlige Verhältnis bringen.
\end{merksatz}
